\documentclass{beamer}

%% Tema y configuraciones
\usetheme{Boadilla}
\usecolortheme{default}
\setbeamertemplate{navigation symbols}{}
\usepackage{xcolor}
\usepackage{tikz}
\usetikzlibrary{arrows.meta, positioning}
\hypersetup{
    colorlinks=true,
    linkcolor=.,      % color de enlaces internos (índices, refs)
    urlcolor=blue,    % color de URLs
    citecolor=.,      % color de referencias bibliográficas
}

%% Helpers (macros y notación)
\newcommand{\E}{\mathbb{E}}
\newcommand{\Var}{\mathrm{Var}}
\newcommand{\Cov}{\mathrm{Cov}}
\newcommand{\1}{\mathbb{I}}
\newcommand{\MSE}{\mathrm{MSE}}
\newcommand{\RSS}{\mathrm{RSS}}
\newcommand{\sign}{\mathrm{sign}}
\DeclareMathOperator*{\argmin}{arg\,min}

%% Encabezado
\title[BDML]{Big Data y Machine Learning \\ para Economía Aplicada}
\subtitle{Week 03: Modelos Lineales de Alta Dimensión --- La Teoría del Lasso}
\author{Eduard F. Martinez Gonzalez, Ph.D.}
\institute[]{Departamento de Economía, Universidad Icesi}
\date{\today}

%%=================================================
%% Begin document
\begin{document}

%%=================================================
%% Primer slide
\begin{frame}
\titlepage
\end{frame}

%%=================================================
\begin{frame}{Previously\ldots}
\small
\begin{itemize}\itemsep5pt
  \item La \textbf{validación cruzada} nos da un estimador honesto del error de generalización --- y con ella giramos perillas.
  \item \textbf{Ridge} $= (\mathbf{X}'\mathbf{X} + \lambda\mathbf{I})^{-1}\mathbf{X}'\mathbf{y}$: encoge proporcionalmente, estabiliza, nunca anula.
  \item \textbf{Lasso}: el valor absoluto produce \textbf{ceros exactos} --- lo vimos con la geometría del diamante y el umbral suave $S_\lambda(z) = \sign(z)(|z|-\lambda)_+$\ldots\ pero ambas fueron \emph{imágenes}, no demostraciones.
\end{itemize}
\pause

\vspace{0.2cm}
\begin{block}{Hoy: abrir el capó (la sesión más matemática del curso)}
\begin{enumerate}\itemsep2pt
  \item ¿Por qué exactamente la alta dimensión rompe a OLS?
  \item ¿De dónde salen \emph{de verdad} los ceros del Lasso? (KKT y subgradientes)
  \item ¿Cómo se calcula? (\emph{coordinate descent}, LARS)
  \item ¿Cuándo recupera las variables verdaderas? (sparsity, irrepresentabilidad)
  \item Y el puente a lo que más nos importa como economistas: usar el Lasso para \textbf{inferencia causal} (\emph{post-double-selection}).
\end{enumerate}
\end{block}
\end{frame}


%%#################################################
%%#################################################
%% PARTE 1: POR QUÉ LA ALTA DIMENSIÓN ROMPE OLS
%%#################################################
%%#################################################
\section{¿Por qué la alta dimensión rompe a OLS?}
\begin{frame}[noframenumbering]
\tableofcontents[currentsection]
\end{frame}

%%=================================================
\begin{frame}{Los tres regímenes de $p$ y $n$}
\small
\begin{enumerate}\itemsep6pt
  \item \textbf{Clásico} ($p$ fijo, $n \to \infty$): el mundo de Econometría I. OLS es consistente, asintóticamente normal; todo funciona. \pause
  \item \textbf{$p$ comparable a $n$}: OLS ``funciona'' pero tiembla --- varianzas enormes, coeficientes inestables. \pause
  \item \textbf{$p \gg n$}: OLS ni siquiera está definido de manera única. \pause
\end{enumerate}

\vspace{0.2cm}
\textbf{¿Y de dónde salen tantos predictores en economía?}
\begin{itemize}\itemsep4pt
  \item Datos administrativos y \emph{scanner data}: miles de características por unidad.
  \item Texto e imágenes convertidos en variables (sesiones 8--9).
  \item \textbf{Flexibilidad que nosotros mismos creamos:} polinomios, \emph{dummies} finas e \textbf{interacciones}. Con 30 variables base y sus interacciones de a pares: $30 + \binom{30}{2} = 465$ regresores. Agreguen cuadrados y triples y $p$ explota.
\end{itemize}
\pause

\begin{block}{El punto}
La alta dimensión no es exótica: aparece apenas uno quiere controlar \emph{flexiblemente} por covariables --- exactamente lo que un economista aplicado hace.
\end{block}
\end{frame}

%%=================================================
\begin{frame}{Qué falla, exactamente (I): $p > n$}
\small
Con $p > n$, la matriz $\mathbf{X}'\mathbf{X}$ ($p \times p$) tiene rango a lo sumo $n < p$:
\begin{itemize}\itemsep5pt
  \item $\mathbf{X}'\mathbf{X}$ es \textbf{singular}: la inversa no existe, las ecuaciones normales tienen \textbf{infinitas soluciones}. \pause
  \item Existen (infinitos) $\beta$ que ajustan los datos \textbf{perfectamente}: $\RSS = 0$, $R^2 = 1$\ldots\ sin que haya ninguna relación real. Es la memorización de la sesión 1 en su forma más pura. \pause
  \item Y no hay grados de libertad para estimar $\sigma^2$: la inferencia clásica muere con el ajuste perfecto.
\end{itemize}
\pause

\vspace{0.2cm}
\begin{block}{Pregunta para la sala}
Si $p > n$ hace que \emph{cualquier} $y$ sea ajustable a la perfección, ¿qué información puede haber en ``ajustar bien''?
\end{block}
\pause
\textcolor{gray}{Ninguna. En alta dimensión, el ajuste dentro de muestra pierde todo su contenido informativo: solo el desempeño \emph{out-of-sample} y la estructura (p.\,ej.\ sparsity) pueden discriminar modelos.}
\end{frame}

%%=================================================
\begin{frame}{Qué falla, exactamente (II): $p \approx n$ y colinealidad}
\small
Aun con $p < n$, si hay predictores muy correlacionados, $\mathbf{X}'\mathbf{X}$ está \emph{cerca} de singular. Descompongamos: sean $d_1 \geq \cdots \geq d_p > 0$ los autovalores de $\mathbf{X}'\mathbf{X}$ con autovectores $v_j$.
\pause

\vspace{0.1cm}
\textbf{La varianza de OLS en la dirección $v_j$:}
\[ \Var\big( v_j'\hat{\beta} \big) \;=\; \sigma^2\, v_j' (\mathbf{X}'\mathbf{X})^{-1} v_j \;=\; \frac{\sigma^2}{d_j} \]
\pause

\begin{itemize}\itemsep5pt
  \item En direcciones con $d_j$ \textbf{pequeño} (combinaciones de predictores casi colineales), la varianza \textbf{explota}: $\sigma^2 / d_j \to \infty$.
  \item Interpretación: los datos casi no varían en esa dirección $\Rightarrow$ casi no informan sobre ella $\Rightarrow$ OLS ``adivina'' con enorme ruido. \pause
  \item Y recuerden la solución de Ridge: reemplazar $d_j$ por $d_j + \lambda$ --- amortiguar exactamente las direcciones enfermas.
\end{itemize}
\pause

\begin{block}{Moraleja}
El enemigo no es $p$ grande \emph{per se}: es la \textbf{información por parámetro}. Alta dimensión $=$ muchos parámetros persiguiendo poca variación independiente.
\end{block}
\end{frame}

%%=================================================
\begin{frame}{La bendición oculta: sparsity}
\small
Si \textbf{todos} los $p$ coeficientes importaran, con $p > n$ el problema sería imposible: más incógnitas que ecuaciones, sin salida.
\pause

\vspace{0.2cm}
\textbf{El supuesto que salva el día:}
\[ \beta^{0} \text{ es \textbf{esparso}: solo } s \ll n \text{ de sus } p \text{ componentes son} \neq 0 \]
(o aproximadamente esparso: pocos grandes, el resto despreciables).
\pause

\begin{itemize}\itemsep5pt
  \item Con sparsity, los ``verdaderos'' grados de libertad son $\sim s$, no $p$: hay esperanza si $n \gg s$, \textbf{aunque} $p \gg n$.
  \item El problema se vuelve de \textbf{búsqueda}: ¿cuáles son las $s$ variables? Búsqueda exacta $= 2^p$ (imposible); el Lasso la vuelve \textbf{convexa}. \pause
  \item ¿Es creíble en economía? A menudo sí como aproximación: de cientos de controles posibles, unos pocos capturan lo esencial. Pero es un \textbf{supuesto}, no un teorema --- y volveremos a él críticamente.
\end{itemize}
\pause

\begin{block}{El programa de hoy}
Entender al Lasso como \emph{la relajación convexa} del problema de búsqueda esparsa --- y estudiar cuándo esa relajación devuelve la respuesta correcta.
\end{block}
\end{frame}


%%#################################################
%%#################################################
%% PARTE 2: EL LASSO COMO PROBLEMA DE OPTIMIZACIÓN
%%#################################################
%%#################################################
\section{El Lasso como problema de optimización}
\begin{frame}[noframenumbering]
\tableofcontents[currentsection]
\end{frame}

%%=================================================
\begin{frame}{Dos formas del mismo problema}
\small
\textbf{Forma penalizada} (la que usamos hasta ahora; convención con $\tfrac{1}{2}$ para derivadas limpias):
\[ \hat{\beta}^{L} \;=\; \argmin_{\beta} \; \frac{1}{2}\,\|\mathbf{y} - \mathbf{X}\beta\|_2^2 \;+\; \lambda\, \|\beta\|_1 \]
\pause
\textbf{Forma restringida:}
\[ \hat{\beta}^{L} \;=\; \argmin_{\beta} \; \frac{1}{2}\,\|\mathbf{y} - \mathbf{X}\beta\|_2^2 \quad \text{sujeto a} \quad \|\beta\|_1 \leq t \]
\pause

\begin{itemize}\itemsep4pt
  \item Son \textbf{equivalentes}: para cada $\lambda > 0$ existe un $t$ (y viceversa) con la misma solución. El puente es el \textbf{Lagrangiano}: $\lambda$ es el multiplicador de la restricción $\|\beta\|_1 \leq t$.
  \item La forma restringida dio la \textbf{geometría} (el diamante); la penalizada es la que se \textbf{computa}. \pause
  \item Ambas son \textbf{convexas}: función objetivo convexa, región factible convexa $\Rightarrow$ todo mínimo local es global, y las condiciones de primer orden \textbf{caracterizan} la solución.
\end{itemize}
\pause

\begin{block}{Por qué esto importa}
Best subset ($\|\beta\|_0 \leq s$) es el problema que \emph{quisiéramos} resolver, pero es combinatorio. El Lasso reemplaza $\|\cdot\|_0$ por $\|\cdot\|_1$: la \textbf{relajación convexa más cercana} que mantiene la esparsidad.
\end{block}
\end{frame}

%%=================================================
\begin{frame}{Repaso exprés: KKT (el Lagrangiano de micro, con desigualdades)}
\small
Para un problema convexo $\min_{\beta} f(\beta)$ s.a.\ $g(\beta) \leq 0$, con $\mathcal{L} = f + \mu g$, la solución $\hat{\beta}$ cumple:
\begin{enumerate}\itemsep4pt
  \item \textbf{Estacionariedad:} $\nabla f(\hat{\beta}) + \mu \nabla g(\hat{\beta}) = 0$.
  \item \textbf{Factibilidad primal:} $g(\hat{\beta}) \leq 0$; \quad \textbf{dual:} $\mu \geq 0$.
  \item \textbf{Holgura complementaria:} $\mu \cdot g(\hat{\beta}) = 0$ (o la restricción no ata, o el precio es positivo).
\end{enumerate}
\pause

\vspace{0.15cm}
\begin{itemize}\itemsep4pt
  \item Es el mismo aparato de la maximización de utilidad con restricción de presupuesto: $\mu$ = precio sombra de la restricción $\|\beta\|_1 \leq t$; aquí ese precio se llama $\lambda$.
  \item En problemas convexos, KKT no es solo necesario: es \textbf{suficiente} --- quien cumple KKT \emph{es} la solución. \pause
\end{itemize}

\vspace{0.1cm}
\begin{block}{El obstáculo}
La estacionariedad pide $\nabla$\ldots\ pero $\|\beta\|_1 = \sum_j |\beta_j|$ \textbf{no es diferenciable} en $\beta_j = 0$ --- que es \emph{justo donde ocurre lo interesante}. Necesitamos generalizar la derivada.
\end{block}
\end{frame}

%%=================================================
\begin{frame}{El subgradiente: derivadas para funciones con esquinas}
\small
\textbf{Definición.} $g$ es un \textbf{subgradiente} de una función convexa $f$ en $x_0$ si
\[ f(x) \;\geq\; f(x_0) + g\,(x - x_0) \qquad \forall x \]
(la recta con pendiente $g$ por $(x_0, f(x_0))$ queda \emph{debajo} de $f$). El conjunto de todos: $\partial f(x_0)$.

\vspace{0.15cm}
\begin{center}
\begin{tikzpicture}[scale=0.62]
  \draw[->] (-2.8,0) -- (2.8,0) node[right, font=\scriptsize] {$\beta_j$};
  \draw[->] (0,-0.4) -- (0,2.6);
  \draw[blue, very thick] (-2.4,2.4) -- (0,0) -- (2.4,2.4);
  \draw[dashed, gray] (-2.2,-1.54) -- (2.2,1.54);   % pendiente 0.7
  \draw[dashed, gray] (-2.2,1.1) -- (2.2,-1.1);     % pendiente -0.5
  \draw[dashed, gray] (-2.4,0) -- (2.4,0);          % pendiente 0
  \node[font=\scriptsize, blue] at (0.0,2.85) {$f(\beta_j) = |\beta_j|$};
  \node[font=\scriptsize, gray, align=left] at (4.6,2.0) {rectas soporte en $0$: \\ cualquier pendiente \\ en $[-1, 1]$};
\end{tikzpicture}
\end{center}
\pause

\textbf{Para el valor absoluto:}
\[ \partial |\beta_j| \;=\; \begin{cases} \{\sign(\beta_j)\} & \text{si } \beta_j \neq 0 \\ [-1, \, 1] & \text{si } \beta_j = 0 \end{cases} \]
\pause
\textbf{Condición de óptimo generalizada:} $\hat{\beta}$ minimiza $f$ convexa $\iff 0 \in \partial f(\hat{\beta})$.
\end{frame}

%%=================================================
\begin{frame}{Las condiciones de optimalidad del Lasso}
\small
Aplicamos $0 \in \partial\big[ \tfrac{1}{2}\|\mathbf{y} - \mathbf{X}\beta\|^2 + \lambda \|\beta\|_1 \big]$. El primer término es diferenciable (gradiente $-\mathbf{X}'(\mathbf{y} - \mathbf{X}\beta)$); el segundo aporta el subgradiente:
\pause
\[ \mathbf{X}_j'(\mathbf{y} - \mathbf{X}\hat{\beta}) \;=\; \lambda\, s_j, \qquad s_j \in \partial|\hat{\beta}_j| \quad \forall j \]
\pause
\textbf{Leyendo por casos:}
\[ \boxed{ \;\begin{aligned}
\hat{\beta}_j \neq 0:& \quad \mathbf{X}_j'(\mathbf{y} - \mathbf{X}\hat{\beta}) = \lambda\, \sign(\hat{\beta}_j) \\
\hat{\beta}_j = 0:& \quad \big| \mathbf{X}_j'(\mathbf{y} - \mathbf{X}\hat{\beta}) \big| \leq \lambda
\end{aligned}\; } \]
\pause

\begin{itemize}\itemsep4pt
  \item \textbf{Interpretación:} $\mathbf{X}_j'(\mathbf{y} - \mathbf{X}\hat{\beta})$ es la covarianza del predictor $j$ con el \textbf{residuo}. Los predictores \textbf{activos} la tienen exactamente en $\pm\lambda$; los \textbf{inactivos}, por debajo en valor absoluto.
  \item $\lambda$ es un \textbf{umbral de mérito}: entras al modelo solo si tu correlación con lo aún no explicado alcanza $\lambda$. \pause
  \item \textbf{Consecuencia inmediata:} si $\lambda \geq \lambda_{\max} \equiv \max_j |\mathbf{X}_j'\mathbf{y}|$, entonces $\hat{\beta} = 0$ cumple las condiciones: \textbf{nadie entra}. El camino de soluciones arranca en $\lambda_{\max}$ --- así construye \texttt{glmnet} su grilla.
\end{itemize}
\end{frame}

%%=================================================
\begin{frame}{Derivación del umbral suave (el átomo del Lasso)}
\small
\textbf{Caso univariado} (o una coordenada con $\mathbf{X}$ estandarizada): sea $z$ el estimador sin penalizar; resolvemos
\[ \hat{\beta} \;=\; \argmin_{\beta} \; \tfrac{1}{2}(z - \beta)^2 + \lambda |\beta| \]
\pause
\textbf{1. Si $\hat{\beta} > 0$:} la condición de primer orden es $\beta - z + \lambda = 0 \Rightarrow \hat{\beta} = z - \lambda$. \\
Coherente con $\hat{\beta} > 0$ solo si $z > \lambda$.
\pause

\textbf{2. Si $\hat{\beta} < 0$:} $\beta - z - \lambda = 0 \Rightarrow \hat{\beta} = z + \lambda$. Requiere $z < -\lambda$.
\pause

\textbf{3. Si $\hat{\beta} = 0$:} necesitamos $0 \in \{-z + \lambda\, s : s \in [-1,1]\}$, es decir $z \in [-\lambda, \lambda]$.
\pause

\vspace{0.15cm}
\textbf{Juntando los tres casos:}
\[ \boxed{ \; \hat{\beta} \;=\; S_{\lambda}(z) \;=\; \sign(z)\,\big(|z| - \lambda\big)_{+} \; } \]
la fórmula que en la sesión 2 mostramos como dibujo, ahora demostrada.
\pause

\begin{block}{Por qué la llamamos ``el átomo''}
Todos los algoritmos modernos del Lasso consisten en aplicar $S_\lambda$ millones de veces, una coordenada a la vez. Eso es lo siguiente.
\end{block}
\end{frame}


%%#################################################
%%#################################################
%% PARTE 3: ALGORITMOS
%%#################################################
%%#################################################
\section{Algoritmos: coordinate descent y LARS}
\begin{frame}[noframenumbering]
\tableofcontents[currentsection]
\end{frame}

%%=================================================
\begin{frame}{Coordinate descent: el algoritmo de \texttt{glmnet}}
\small
\textbf{Idea:} optimizar \textbf{una coordenada a la vez}, dejando las demás fijas --- porque el subproblema de una coordenada \emph{es} el umbral suave que acabamos de derivar.
\pause

\vspace{0.15cm}
\textbf{Algoritmo} (con predictores estandarizados, objetivo $\tfrac{1}{2n}\|\mathbf{y}-\mathbf{X}\beta\|^2 + \lambda\|\beta\|_1$):
\begin{enumerate}\itemsep3pt
  \item Inicializar $\hat{\beta} = 0$.
  \item Repetir hasta converger, para $j = 1, \ldots, p$:
  \begin{itemize}\itemsep2pt
    \item Residuo parcial: $\; r^{(j)} = \mathbf{y} - \sum_{k \neq j} \mathbf{X}_k \hat{\beta}_k$
    \item Actualizar: $\; \hat{\beta}_j \leftarrow S_{\lambda}\!\big( \tfrac{1}{n}\, \mathbf{X}_j' r^{(j)} \big)$
  \end{itemize}
\end{enumerate}
\pause

\textbf{¿Por qué converge?} El objetivo es convexo y la parte no diferenciable ($\|\beta\|_1$) es \textbf{separable} por coordenadas --- la condición exacta bajo la cual el descenso coordenado alcanza el óptimo global (Tseng, 2001).
\pause

\vspace{0.1cm}
\textbf{Los trucos de \texttt{glmnet}} (Friedman, Hastie \& Tibshirani, 2010):
\begin{itemize}\itemsep2pt
  \item \emph{Warm starts}: resolver el camino completo de $\lambda_{\max}$ hacia abajo, usando cada solución como arranque de la siguiente.
  \item \emph{Active set}: iterar solo sobre las variables activas y chequear el resto de vez en cuando (por KKT, la mayoría permanece en cero).
\end{itemize}
\end{frame}

%%=================================================
\begin{frame}{LARS y la forma del camino}
\small
\textbf{Least Angle Regression} (Efron, Hastie, Johnstone \& Tibshirani, 2004): la alternativa geométrica.
\begin{itemize}\itemsep4pt
  \item Arrancar en $\beta = 0$; avanzar en la dirección del predictor más correlacionado con el residuo.
  \item Cuando otro predictor \textbf{alcanza} esa correlación, avanzar en la dirección \textbf{equiangular} entre ambos; y así sucesivamente.
  \item Con una pequeña modificación, produce \textbf{exactamente} el camino del Lasso, al costo de un solo OLS. \pause
\end{itemize}

\vspace{0.1cm}
\begin{center}
\begin{tikzpicture}[scale=0.66]
  \draw[->] (0,-1.5) -- (0,2.5) node[above, font=\scriptsize] {$\hat{\beta}_j$};
  \draw[->] (0,0) -- (7.9,0);
  \node[font=\scriptsize] at (3.9,-1.95) {$\|\hat{\beta}\|_1$ crece $\rightarrow$ \; ($\lambda$ decrece)};
  \draw[blue, very thick] (0.8,0) -- (2.6,0.85) -- (4.4,1.5) -- (6.2,1.9) -- (7.4,2.15);
  \draw[red, very thick] (2.6,0) -- (4.4,0.7) -- (6.2,1.15) -- (7.4,1.4);
  \draw[teal, very thick] (4.4,0) -- (6.2,-0.6) -- (7.4,-0.95);
  \draw[orange, very thick] (6.2,0) -- (7.4,0.3);
  \foreach \x in {0.8, 2.6, 4.4, 6.2} \draw[gray, dashed] (\x,-1.3) -- (\x,2.2);
\end{tikzpicture}
\end{center}
\pause

\begin{itemize}\itemsep3pt
  \item \textbf{El camino del Lasso es lineal por tramos:} los quiebres (líneas grises) son entradas (o salidas) de variables del conjunto activo.
  \item Hoy LARS es sobre todo una \textbf{fuente de intuición}; en producción domina \emph{coordinate descent} (escala mejor y cubre logística, Elastic Net, etc.).
\end{itemize}
\end{frame}


%%#################################################
%%#################################################
%% PARTE 4: CUÁNDO EL LASSO RECUPERA LA VERDAD
%%#################################################
%%#################################################
\section{¿Cuándo el Lasso recupera la verdad?}
\begin{frame}[noframenumbering]
\tableofcontents[currentsection]
\end{frame}

%%=================================================
\begin{frame}{Dos preguntas que suenan igual y no lo son}
\small
Supongamos que el mundo es $\mathbf{y} = \mathbf{X}\beta^0 + \varepsilon$ con $\beta^0$ esparso (soporte $S$, $|S| = s$).
\pause

\vspace{0.2cm}
\begin{enumerate}\itemsep6pt
  \item \textbf{¿Predice bien?} ¿Es pequeño el error $\tfrac{1}{n}\|\mathbf{X}(\hat{\beta} - \beta^0)\|_2^2$? \pause
  \item \textbf{¿Selecciona bien?} ¿Recupera el soporte: $\{j : \hat{\beta}_j \neq 0\} = S$? \pause
\end{enumerate}

\vspace{0.2cm}
\begin{itemize}\itemsep5pt
  \item La (2) es \textbf{mucho} más exigente que la (1): se puede predecir de maravilla con las variables ``equivocadas'' (proxies correlacionados).
  \item Y son relevantes para usos distintos: para \textbf{pronosticar}, basta (1); para \textbf{interpretar} qué variables importan --- o elegir \emph{controles} --- quisiéramos (2).
\end{itemize}
\pause

\begin{block}{Lo que sigue}
Los dos resultados centrales de la teoría, presentados con su \textbf{intuición} (las demostraciones: Wainwright, caps.\ 7--8, para quien quiera el detalle).
\end{block}
\end{frame}

%%=================================================
\begin{frame}{Resultado 1: el precio de la alta dimensión es solo $\log p$}
\small
\textbf{Teorema (informal).} Con errores razonables, $\lambda \asymp \sigma\sqrt{\log(p)/n}$, y un diseño no demasiado degenerado en las direcciones esparsas (\emph{restricted eigenvalue}),
\[ \frac{1}{n}\, \big\| \mathbf{X}(\hat{\beta}^L - \beta^0) \big\|_2^2 \;\lesssim\; \sigma^2\, \frac{s \, \log p}{n} \]
\pause

\textbf{Cómo leerlo:}
\begin{itemize}\itemsep5pt
  \item Si conociéramos de antemano las $s$ variables correctas, OLS sobre ellas daría error $\sim \sigma^2 s / n$. \pause
  \item El Lasso, \textbf{sin saber cuáles son} entre $p$ candidatas, paga solo un factor $\log p$ adicional: el costo de la búsqueda es \textbf{logarítmico}. \pause
  \item Por eso $p$ puede ser enorme: con $p = 10^6$, $\log p \approx 14$. Lo que de verdad limita es $s$ (cuánta señal hay) y $n$ (cuántos datos).
\end{itemize}
\pause

\begin{block}{La condición del diseño, en palabras}
Se necesita que la colinealidad no destruya la información \emph{en las direcciones esparsas relevantes}. Es mucho más débil que pedir $\mathbf{X}'\mathbf{X}$ invertible --- y por eso el Lasso vive donde OLS muere.
\end{block}
\end{frame}

%%=================================================
\begin{frame}{Resultado 2: recuperar el soporte exige mucho más}
\small
Para que $\{j : \hat{\beta}_j \neq 0\} = S$ con alta probabilidad se necesita, además:
\pause
\begin{enumerate}\itemsep5pt
  \item \textbf{Señales no minúsculas} (\emph{beta-min}): $\min_{j \in S} |\beta^0_j| \gtrsim \lambda$. Coeficientes verdaderos por debajo del umbral de ruido son indistinguibles de cero --- nadie puede encontrarlos. \pause
  \item \textbf{Condición de irrepresentabilidad} (Zhao \& Yu, 2006):
  \[ \big\| \mathbf{X}_{S^c}'\mathbf{X}_S \big( \mathbf{X}_S'\mathbf{X}_S \big)^{-1} \sign(\beta^0_S) \big\|_\infty \;<\; 1 \]
  En palabras: \textbf{ningún predictor irrelevante puede ser demasiado bien explicado (regresado) por los relevantes}. \pause
\end{enumerate}

\vspace{0.15cm}
\begin{itemize}\itemsep4pt
  \item Es una condición \textbf{fuerte} y --- lo incómodo --- \textbf{no verificable} (depende del soporte $S$, que no conocemos).
  \item Si falla, el Lasso incluye proxies equivocados \emph{sin importar cuántos datos haya}: no es un problema de $n$, es estructural.
\end{itemize}
\end{frame}

%%=================================================
\begin{frame}{El modo de falla típico: el proxy correlacionado}
\small
\begin{center}
\begin{tikzpicture}[
    every node/.style={circle, draw, minimum size=0.9cm, font=\scriptsize},
    arrow/.style={-{Stealth[length=2.5mm]}, thick}
]
  \node (X1) {$X_1$};
  \node (X2) [right=2.6cm of X1] {$X_2$};
  \node (Y) [below right=0.95cm and 0.6cm of X1] {$Y$};
  \draw[arrow, blue, thick] (X1) -- node[draw=none, circle, left, font=\scriptsize\color{blue}]{$\beta_1 \neq 0$} (Y);
  \draw[{Stealth[length=2.5mm]}-{Stealth[length=2.5mm]}, dashed, red, thick] (X1) -- node[draw=none, circle, above, font=\scriptsize\color{red}]{$\rho = 0.95$} (X2);
\end{tikzpicture}
\end{center}

\begin{itemize}\itemsep2pt
  \item $X_1$ es la variable verdadera; $X_2$ es un proxy muy correlacionado que \textbf{no} afecta a $Y$.
  \item Para el ojo del Lasso, ambos predicen casi igual: con ruido de por medio, puede quedarse con $X_2$ (o repartir, o alternar entre muestras). \pause
  \item Para \textbf{predecir}, da lo mismo; para \textbf{interpretar}, es un desastre silencioso: el modelo reporta con confianza la variable equivocada.
\end{itemize}
\pause

\begin{block}{La frase para tatuarse}
El Lasso selecciona \textbf{predictores}, no \textbf{causas}.
\end{block}
\end{frame}

%%=================================================
\begin{frame}{¿Y los métodos secuenciales? (forward/backward vs.\ Lasso)}
\small
\begin{center}
\footnotesize
\setlength{\tabcolsep}{4pt}
\begin{tabular}{lll}
\hline
 & Forward/backward stepwise & Lasso \\
\hline
Naturaleza & búsqueda \emph{greedy} discreta & optimización convexa global \\
Solución & depende del camino & única (diseño general) \\
Estabilidad & baja: todo o nada & mayor: encogimiento continuo \\
Garantías teóricas & débiles & tasas y condiciones claras \\
Costo & $O(p^2)$ ajustes & camino completo, rápido \\
\hline
\end{tabular}
\end{center}
\pause

\vspace{0.25cm}
\begin{itemize}\itemsep5pt
  \item El stepwise \emph{hereda} todos los problemas de selección del Lasso (proxies, inestabilidad) \textbf{sin} sus garantías.
  \item Y ambos comparten el pecado capital para nosotros: \textbf{la inferencia clásica después de seleccionar es inválida} --- los errores estándar fingen que el modelo estaba decidido antes de ver los datos.
\end{itemize}
\pause

\vspace{0.1cm}
Ese pecado es exactamente el tema del cierre de hoy.
\end{frame}


%%#################################################
%%#################################################
%% PARTE 5: POST-DOUBLE-SELECTION
%%#################################################
%%#################################################
\section{Del Lasso a la inferencia causal: post-double-selection}
\begin{frame}[noframenumbering]
\tableofcontents[currentsection]
\end{frame}

%%=================================================
\begin{frame}{El problema que nos importa como economistas}
\small
\textbf{Setup:} queremos el efecto de un tratamiento $d$ (no aleatorizado) sobre $y$, controlando \emph{flexiblemente} por confusores observables:
\[ y_i \;=\; \alpha\, d_i \;+\; x_i'\beta \;+\; \varepsilon_i, \qquad p = \dim(x) \text{ grande} \]
\begin{itemize}\itemsep2pt
  \item $x$: cientos de controles candidatos (variables base $+$ polinomios $+$ interacciones).
  \item El objetivo es $\alpha$ \textbf{con un intervalo de confianza válido} --- no predecir $y$.
\end{itemize}
\pause

\textbf{La tentación (single selection):}
\begin{enumerate}\itemsep2pt
  \item Lasso de $y$ sobre $(d, x)$ sin penalizar $d$ $\Rightarrow$ controles seleccionados $\hat{S}$.
  \item OLS de $y$ sobre $d$ y $x_{\hat{S}}$; reportar el error estándar de siempre.
\end{enumerate}
\pause

\begin{block}{Pregunta para la sala}
¿Dónde está el error?
\end{block}
\pause
\textcolor{gray}{El Lasso selecciona por \textbf{capacidad de predecir $y$} --- pero un confusor peligroso no necesita eso: le basta estar \textbf{muy correlacionado con $d$}. A ese, el Lasso lo bota, y su omisión sesga $\hat{\alpha}$.}
\end{frame}

%%=================================================
\begin{frame}{El fallo, en el lenguaje de su curso de causalidad}
\small
\begin{center}
\begin{tikzpicture}[
    every node/.style={circle, draw, minimum size=1.05cm, font=\scriptsize},
    arrow/.style={-{Stealth[length=2.5mm]}, thick}
]
  \node (Z) {$x_k$};
  \node (D) [below left=1.3cm and 1.9cm of Z] {$d$};
  \node (Y) [below right=1.3cm and 1.9cm of Z] {$y$};
  \draw[arrow, red] (Z) -- node[draw=none, circle, left, font=\scriptsize\color{red}]{fuerte} (D);
  \draw[arrow, red] (Z) -- node[draw=none, circle, right, font=\scriptsize\color{red}]{débil (pero $\neq 0$)} (Y);
  \draw[arrow, blue, thick] (D) -- node[draw=none, circle, below, font=\scriptsize\color{blue}]{$\alpha$} (Y);
\end{tikzpicture}
\end{center}
\pause

\begin{itemize}\itemsep4pt
  \item $x_k$: confusor con efecto \textbf{débil sobre $y$} pero \textbf{fuertemente correlacionado con $d$}.
  \item El Lasso de $y$ sobre $(d,x)$ lo descarta ($d$ ya ``absorbe'' su poca señal sobre $y$) $\Rightarrow$ backdoor abierto $\Rightarrow$ \textbf{sesgo de variable omitida} en $\hat{\alpha}$: el $\delta \cdot \phi$ de su curso de inferencia causal, resucitado por un algoritmo. \pause
  \item Y el problema es sistemático: la selección se equivoca \emph{precisamente} en los regresores de coeficiente moderado --- los errores de selección son inevitables en muestras finitas.
\end{itemize}
\pause

\begin{block}{La solución no es seleccionar mejor}
Es seleccionar \textbf{dos veces}, para que los errores de selección no importen.
\end{block}
\end{frame}

%%=================================================
\begin{frame}{Post-double-selection (Belloni, Chernozhukov \& Hansen, 2014)}
\small
\textbf{El algoritmo:}
\begin{enumerate}\itemsep5pt
  \item \textbf{Lasso de $y$ sobre $x$:} controles que predicen el \emph{resultado} $\Rightarrow \hat{S}_y$.
  \item \textbf{Lasso de $d$ sobre $x$:} controles que predicen el \emph{tratamiento} $\Rightarrow \hat{S}_d$.
  \item \textbf{OLS de $y$ sobre $d$ y $x_{\hat{S}_y \cup \hat{S}_d}$}; reportar los errores estándar robustos de siempre.
\end{enumerate}
\pause

\vspace{0.15cm}
\textbf{Teorema (informal):} bajo sparsity aproximada de ambas ecuaciones, $\hat{\alpha}$ es $\sqrt{n}$-consistente y asintóticamente normal; los intervalos de confianza usuales son \textbf{uniformemente válidos} --- incluso aunque la selección cometa errores.
\pause

\vspace{0.15cm}
\textbf{Por qué funciona (la intuición que hay que llevarse):}
\begin{itemize}\itemsep4pt
  \item Para sesgar $\hat{\alpha}$, una variable omitida debe correlacionar con $y$ \textbf{y} con $d$ (el $\delta \cdot \phi$).
  \item PDS solo omite variables débiles \textbf{en ambas} ecuaciones $\Rightarrow$ el sesgo residual es un producto de dos términos pequeños: de \textbf{segundo orden}, despreciable frente a $1/\sqrt{n}$.
  \item Ese diseño ``inmunizado'' a errores de selección tiene nombre propio --- \textbf{ortogonalidad de Neyman} --- y es la semilla del \emph{double machine learning} (sesión 6).
\end{itemize}
\end{frame}

%%=================================================
\begin{frame}{PDS en la práctica}
\small
\begin{itemize}\itemsep6pt
  \item \textbf{En R:} paquete \texttt{hdm} (Chernozhukov, Hansen \& Spindler): \texttt{rlassoEffect()} implementa PDS con una línea; reporta $\hat{\alpha}$, SE e IC válidos. \pause
  \item \textbf{Detalle fino:} para inferencia, $\lambda$ \textbf{no} se elige por CV sino con la fórmula \emph{plug-in} teórica ($\lambda \asymp \sigma\sqrt{n\log p}$, con ajustes): el CV optimiza predicción y tiende a dejar entrar demasiadas variables --- otro recordatorio de que predicción e inferencia son juegos distintos. \pause
  \item \textbf{Dónde lo van a ver:} regresiones de crecimiento (¿qué controles entre docenas?), evaluación con muchos confusores administrativos, efectos de políticas con controles flexibles (BCH lo ilustran re-examinando la relación aborto--crimen y leyes de armas). \pause
\end{itemize}

\vspace{0.1cm}
\begin{block}{Regla práctica para sus proyectos}
¿Muchos controles y un $\alpha$ que interpretar? \textbf{Nunca} un solo Lasso de $y$: siempre doble selección (o, en general, un procedimiento ortogonalizado).
\end{block}
\end{frame}

%%=================================================
\begin{frame}{Recap}
\small
\begin{enumerate}\itemsep6pt
  \item La alta dimensión rompe a OLS por \textbf{información por parámetro}: $\Var(v_j'\hat{\beta}) = \sigma^2/d_j$ explota en direcciones casi colineales; con $p > n$, el ajuste perfecto pierde todo contenido.
  \item La salida es un supuesto de estructura: \textbf{sparsity}. El Lasso es la \textbf{relajación convexa} de la búsqueda esparsa.
  \item \textbf{KKT $+$ subgradiente} caracterizan la solución: activos en $|\mathbf{X}_j'r| = \lambda$, inactivos por debajo; de ahí el umbral suave $S_\lambda$, demostrado por casos.
  \item \emph{Coordinate descent} $=$ aplicar $S_\lambda$ coordenada por coordenada; el camino es lineal por tramos (LARS).
  \item \textbf{Predicción} cuesta solo $\log p$; \textbf{selección} exige irrepresentabilidad (fuerte, incomprobable): el Lasso selecciona \emph{predictores}, no \emph{causas}.
  \item Para inferencia sobre un tratamiento: \textbf{post-double-selection} --- seleccionar con $y$ \emph{y} con $d$, y la validez sobrevive a los errores de selección (ortogonalidad de Neyman).
\end{enumerate}
\end{frame}

%%=================================================
\begin{frame}
\vspace{2.0cm}
\Large{Preparación para la próxima clase\ldots}
\end{frame}

%%#################################################
%%#################################################
%% PARTE 6: PRÓXIMA CLASE
%%#################################################
%%#################################################

%%=================================================
\begin{frame}{Para aprovechar al máximo la próxima sesión}
\small
\textbf{Lecturas obligatorias de esta semana:}
\begin{itemize}\itemsep4pt
  \item Belloni, Chernozhukov \& Hansen (2014), \emph{High-Dimensional Methods and Inference on Structural and Treatment Effects}, \emph{JEP}. \\ \textcolor{gray}{Guía: la versión amable de la sesión de hoy; léanla completa.}
  \item Hastie, Tibshirani \& Wainwright (2015), \emph{Statistical Learning with Sparsity}, caps.\ 1--2. \\ \textcolor{gray}{Guía: formalizan lo visto; el cap.\ 11 es opcional para los valientes.}
\end{itemize}

\textbf{Recomendadas:} Efron et al.\ (2004, LARS); Wainwright (2019), caps.\ 7--8 (nivel doctorado); Zhao \& Yu (2006).
\pause

\vspace{0.2cm}
\textbf{Próxima sesión --- cambio total de familia:} dejamos los modelos lineales. Árboles de decisión, Random Forest y \emph{boosting}: los métodos que suelen ganar con datos tabulares. Lean ISLR cap.\ 8 \textbf{antes} de clase.
\pause

\vspace{0.2cm}
\textbf{Aplicación de esta semana en R:} precios de vivienda con muchas covariables --- Lasso y Elastic Net con \texttt{glmnet}, y PDS con \texttt{hdm}.
\end{frame}

%%=================================================
%% Referencias
\begin{frame}{Referencias esenciales}
\footnotesize
\begin{itemize}\itemsep2pt
  \item Tibshirani, R. (1996). Regression Shrinkage and Selection via the Lasso. \emph{JRSS-B}, 58(1), 267--288.
  \item Efron, B., Hastie, T., Johnstone, I., \& Tibshirani, R. (2004). Least Angle Regression. \emph{Annals of Statistics}, 32(2), 407--499.
  \item Zhao, P., \& Yu, B. (2006). On Model Selection Consistency of Lasso. \emph{JMLR}, 7, 2541--2563.
  \item Friedman, J., Hastie, T., \& Tibshirani, R. (2010). Regularization Paths for GLMs via Coordinate Descent. \emph{J.\ of Statistical Software}, 33(1).
  \item Hastie, T., Tibshirani, R., \& Wainwright, M. (2015). \emph{Statistical Learning with Sparsity}. CRC Press.
  \item Wainwright, M. (2019). \emph{High-Dimensional Statistics: A Non-Asymptotic Viewpoint}. Cambridge UP. Caps.\ 7--8.
  \item Leeb, H., \& Pötscher, B. (2005). Model Selection and Inference: Facts and Fiction. \emph{Econometric Theory}, 21(1), 21--59.
  \item Belloni, A., Chernozhukov, V., \& Hansen, C. (2014). Inference on Treatment Effects after Selection among High-Dimensional Controls. \emph{Review of Economic Studies}, 81(2), 608--650.
  \item Belloni, A., Chernozhukov, V., \& Hansen, C. (2014). High-Dimensional Methods and Inference on Structural and Treatment Effects. \emph{JEP}, 28(2), 29--50.
  \item Chernozhukov, V., Hansen, C., \& Spindler, M. (2016). hdm: High-Dimensional Metrics. \emph{R Journal}, 8(2).
\end{itemize}
\normalsize
\end{frame}

%%=================================================
%% End document
\end{document}
